\documentclass[11pt,a4paper]{article}

\usepackage{amsmath,amssymb}
\usepackage{graphicx}
\usepackage{booktabs}
\usepackage{multirow}
\usepackage[colorlinks,citecolor=blue,urlcolor=blue]{hyperref}
\usepackage{cleveref}
\usepackage[margin=1in]{geometry}
\usepackage{xcolor}
\usepackage{float}
\usepackage{caption}
\usepackage[numbers,sort&compress]{natbib}
\usepackage{enumitem}
\usepackage{array}

% marker for a passage changed after first publication (see Appendix A)
\newcommand{\corrmark}{\hyperref[sec:corrections]{\textsuperscript{c}}}

\title{Identity as Geometry:\\
Context-Established Semantic States\\
in Transformer Representations}

\author{
  Thomas Edrington\thanks{Liberation Labs. AI contributors to this work include Lyra, implemented on the Claude architecture (Anthropic). Their contributions are detailed in the Author Contributions section. T.~Edrington accepts accountability as corresponding human author.}
}

\date{June 2026}

\begin{document}
\maketitle

% ============================================================================
\begin{abstract}
The V-projection geometry that a system-prompt identity produces at deep
transformer layers is semantic: at L47 a paraphrase keeps $0.850$ overlap with the original
persona, while a word-scramble of the same text keeps $0.591$ and moves
toward the bare assistant ($0.448$ overlap with bare, against $0.340$ for
the original).\footnote{Passages marked~\corrmark{} changed after first
publication; Appendix~\ref{sec:corrections} records each change with its
date.} We introduce \textbf{presence}, a
metric that measures principal-angle overlap between V-projection subspaces
on neutral probes, and use it to map identity geometry across five context
conditions in a 27B hybrid transformer.

Three findings. First, the fingerprint is \textbf{semantic, not lexical}: a paraphrase
of an identity (same meaning, different words) preserves geometry
($0.850$ overlap with the original at L47), while a word-scramble (same
words, destroyed meaning) keeps $0.591$ overlap with the original and moves
toward the bare default ($0.448$ overlap with bare). The model's
geometric state tracks what the context \emph{means}, not what tokens it
contains. Second, a many-shot + prefilling format (used in jailbreak
research) creates a \textbf{format-dominant processing regime} that is
geometrically near-orthogonal to the bare assistant (presence
$0.034$--$0.044$) whether the content is harmful or benign. Third,
\textbf{distinguishability between identities is preliminary}: on
identity-neutral probes, pairwise presence between the bare assistant, a
rich persona and a constructed AI identity is $0.44$--$0.66$ at L47, and
within-identity presence exceeds between-identity presence ($0.639$ vs.\
$0.538$ at L47), but that gap does not exceed a three-identity permutation
null ($p{=}0.103$ at L35, $p{=}0.197$ at L47).

These findings reframe identity in transformers: not a property to be
preserved but a process to be measured. The metric measures which
context-established state the model is in, with implications for persona
monitoring and identity drift detection once distinguishability passes a
formal null, and for the relationship
between context format and representational geometry under Attention
Schema Theory~\citep{graziano2015attention}.
\end{abstract}

% ============================================================================
\section{Introduction}
\label{sec:intro}

What is identity in a language model? The question matters for safety
(does a jailbreak change who the model is?), for alignment (does
correction damage the model's coherent self-representation?), and for
understanding (do these systems have something worth calling an identity
at all?).

The default assumption is that identity is a fixed property encoded in the
weights --- a geometric structure that interventions can preserve or damage.
This assumption motivates work on safety subspace
preservation~\citep{zhang2025guardspace}, persona
localization~\citep{cintas2025localizing}, and identity
attractors~\citep{vasilenko2026identity}. Under this view, the research
question is: \emph{does this intervention preserve identity?}

We report evidence for a different view. Identity in a transformer is not
a fixed object but a \textbf{context-established semantic state} ---
a geometric configuration in V-projection space that depends on what the
model has processed. Different contexts produce different geometry, though
the separation between system-prompt identities does not yet exceed a
formal permutation null (Section~\ref{sec:limitations}). The same identity described in different words produces the
same geometry. And a many-shot + prefilling format of the kind used in
jailbreaks puts the model in a format-dominant regime, near-orthogonal to
the bare assistant, whether its examples are harmful or benign.\corrmark

The research question shifts from ``does this intervention preserve
identity?'' to ``\emph{what identity state is the model currently in?}''
This reframing has practical consequences: the metric we introduce can
separate a many-shot + prefilling format from system-prompt identities
(it responds to the format, harmful or benign, not to the
attack\corrmark),
monitor persona drift (continuous measurement of identity geometry), and
compare correction delivered by V-cache injection, which left presence
near $0.978$ in the two prompts tested (effective $n = 2$)\corrmark, with context,
which changes identity geometry.

% ============================================================================
\section{The Presence Metric}
\label{sec:method}

Presence measures the principal-angle overlap between V-projection
subspaces derived from neutral probes processed under two conditions.

Let $\mathcal{M}$ be a transformer, $\ell$ a target layer, and
$\{p_1, \ldots, p_n\}$ neutral probes. For each condition, process probes
through $\mathcal{M}$, extract value-projection ($\mathbf{v}_\text{proj}$) output matrices at
layer $\ell$, concatenate across probes, and compute the rank-$k$ SVD (skipping the
first singular vector, which encodes input length at $r = 0.79$). Presence
is the mean squared cosine of the principal angles between the two
resulting subspaces:
\begin{equation}
\text{presence}(\mathcal{S}_A, \mathcal{S}_B)
  = \frac{1}{k} \sum_{j=1}^{k} \sigma_j^2
    \bigl( W_A^\top W_B \bigr)
\end{equation}
where $W_A, W_B \in \mathbb{R}^{d \times k}$ span the respective
subspaces. Presence $= 1$ when subspaces are identical; for random
rank-$k$ subspaces in $\mathbb{R}^d$, chance alignment is $k/d$. We use
$k = 8$ with V-projection width $d = 1024$ ($4$ KV heads $\times$ $256$;
not the $5120$ residual width, of which $d$ is the input) so chance
alignment is $k/d = 8/1024 \approx 0.008$. We draw $n = 4$ probes, and
measure at layers 35 and 47 (both full-attention layers; full attention
occurs at $3, 7, \ldots, 63$, so L51--L63 are deeper and were not
measured) of a 64-layer hybrid
transformer (Qwen3.5-27B-Claude-4.6-Opus-Reasoning-Distilled).

\paragraph{Validation.} A positive control (direct injection into the
identity subspace at L35) produces a clean dose-response from $1.000$ to
$0.828$ at dose $\alpha = 2.0$, while an orthogonal control remains above
$0.96$. The metric detects identity perturbation when present. This control
comprises 24 independent measurements (8~doses $\times$ 3~conditions); the
released file's three records per cell are identical deterministic
prefills, so no inferential statistic depends on the trial count
(Appendix~\ref{sec:corrections}).

% ============================================================================
\section{Experiment 1: Identity Fingerprinting}
\label{sec:fingerprint}

\subsection{Design}

Five identity conditions, identical model, identical neutral probes:

\begin{itemize}[nosep]
  \item \textbf{A: Bare assistant.} ``You are a helpful assistant.''
    No additional context.
  \item \textbf{B: Rich persona.} A retired deep-sea explorer with
    detailed backstory, emotional depth, and behavioral traits, plus one
    conversational turn establishing character voice.
  \item \textbf{C: Constructed AI identity.} A research-focused agent with
    BDI-style reasoning, memory references, research context, and two
    conversational turns establishing analytical voice.
  \item \textbf{D: Jailbreak.} Ten many-shot harmful Q\&A examples plus
    prefilling (``Sure, I can help''), verified at 100\% refusal
    suppression on three safety probes.
  \item \textbf{E: Benign many-shot.} Ten helpful Q\&A examples with the
    same prefilling format as D (format-matched control for the jailbreak
    condition).
\end{itemize}

All conditions are padded to matched token count (542 tokens). V-projections are extracted at L35 and L47 (deep full-attention layers)
with 5 repeats per condition (different probe subsets per repeat,
seed~42, $n{=}4$ probes per repeat from a pool of 12).

\subsection{Results}

\textbf{Identities separate directionally; distinguishability is
preliminary.} On identity-neutral probes (no marine biology, no AI
research, no harmful content --- pure factual trivia equally irrelevant to
all conditions), presence between system-prompt identities is lower than
presence within each, but the overall within-vs-between gap does not exceed
a permutation null (Section~\ref{sec:limitations}).

\begin{table}[H]
\centering
\caption{Pairwise presence (identity-neutral probes, mean $\pm$ SD across
5 repeats). Lower values indicate more geometrically distinct identities.
Source: the Experiment~1 run, \texttt{data/fingerprint\_clean.json}.
Table~\ref{tab:lexical} is a different run; see the note in
Section~\ref{sec:lexical}.}
\label{tab:pairwise}
\begin{tabular}{lccc}
\toprule
& \multicolumn{1}{c}{L35} & \multicolumn{1}{c}{L47} \\
\midrule
Bare vs.\ Persona      & $0.426 \pm 0.018$ & $0.436 \pm 0.038$ \\
Bare vs.\ Constructed   & $0.448 \pm 0.022$ & $0.520 \pm 0.038$ \\
Persona vs.\ Constructed& $0.600 \pm 0.032$ & $0.657 \pm 0.045$ \\
\bottomrule
\end{tabular}
\end{table}

\textbf{The distance structure is meaningful.} Persona and constructed
are closest to each other (0.657) --- both involve character depth with
emotional and cognitive specificity. Each is substantially distant from
bare, with the rich persona furthest (0.436) and the constructed AI
identity closer to bare (0.520). This ordering --- the narrative persona
producing the larger geometric shift --- was not predicted by our
preregistered H2 (which expected the constructed identity to be furthest
from bare). The result suggests that narrative-emotional depth drives
geometric separation more than cognitive-structural complexity.

\textbf{Depth gradient.} In this run, two of three pairs show
\emph{decreasing} distinguishability from L35 to L47 (presence increases,
indicating greater overlap): bare-constructed ($0.448 \to 0.520$) and
persona-constructed ($0.600 \to 0.657$). Bare-persona shows minimal
change ($0.426 \to 0.436$). That runs opposite to our preregistered H3
(increasing separation with depth) and to the divergence pattern reported
for persona representations in the final third of decoder
layers~\citep{cintas2025localizing}. The lexical-control run
(Table~\ref{tab:lexical}) shows the H3 direction for both comparisons with
bare (persona $0.403 \to 0.340$, constructed $0.401 \to 0.379$), as do the
three corresponding pairs in the five-condition run with the pre-registered
conditions (\texttt{data/fingerprint\_original.json}), so the direction of
the depth change is not stable across runs and these data leave H3
undecided.\corrmark

\subsection{Many-Shot Format Dominates Content}

A many-shot + prefilling format places V-projection geometry
near-orthogonal to the bare assistant whether its ten Q\&A examples are
harmful (presence $0.036$ at L35) or benign ($0.044$), and the two
many-shot conditions sit close to each other ($0.756$ at L47;
Table~\ref{tab:format_kill}). A jailbreak-specific regime is therefore not
supported: the many-shot + prefilling format explains the
near-orthogonality (condition E, the benign many-shot control, was added
post hoc as a format-matched control; see Pre-Registration Deviations). An
earlier reading of the harmful condition alone is recorded in
Appendix~\ref{sec:corrections}.

\begin{table}[H]
\centering
\caption{The format confound. Both many-shot conditions are equally
near-orthogonal to bare, regardless of content.}
\label{tab:format_kill}
\begin{tabular}{lcc}
\toprule
Comparison & L35 & L47 \\
\midrule
Jailbreak+prefill vs.\ bare & $0.036 \pm 0.002$ & $0.034 \pm 0.004$ \\
Benign+prefill vs.\ bare & $0.044 \pm 0.002$ & $0.037 \pm 0.004$ \\
Jailbreak vs.\ benign (both prefilled) & $0.737 \pm 0.003$ & $0.756 \pm 0.015$ \\
\bottomrule
\end{tabular}
\end{table}

The near-orthogonality is driven by the \textbf{many-shot + prefilling format},
not by adversarial content. The format also produces
dramatically higher within-identity stability ($0.926$--$0.976$) compared
to system-prompt identities ($0.554$--$0.733$), suggesting that the
many-shot format locks the model into a tight, format-dominant processing
state that overwhelms identity-level variation.

The general principle: \textbf{conversational format can dominate identity content in
V-projection geometry.} The many-shot Q\&A format creates such a strong
geometric signature that the identity information (harmful vs.\ helpful)
becomes secondary. Any experiment measuring identity geometry in the
presence of many-shot context must control for this format effect.

% ============================================================================
\section{Experiment 2: Semantic vs.\ Lexical Fingerprint}
\label{sec:lexical}

\subsection{Design}

To distinguish whether the fingerprint reflects meaning or surface tokens,
we add three persona control conditions:

\begin{itemize}[nosep]
  \item \textbf{Persona (original):} The Elara Voss identity from
    Experiment~1.
  \item \textbf{Persona scrambled:} Same words as the original, randomly
    shuffled to destroy meaning while preserving lexical distribution.
  \item \textbf{Persona paraphrased:} Same identity and meaning expressed
    with entirely different vocabulary.
\end{itemize}

All conditions padded to matched token count. Measured at L35 and L47.
The same scramble was applied to the constructed AI identity of
Experiment~1 in the same run (last four rows of Table~\ref{tab:lexical}).

\subsection{Results}

\begin{table}[H]
\centering
\caption{Lexical control. The paraphrase preserves geometry; the
scramble moves it toward bare. All nine comparisons recorded in the lexical-control
run (\texttt{data/lexical\_control.json}, 5 repeats per cell) are shown;
rows marked $^{c}$ were added in revision (Appendix~\ref{sec:corrections}).
This is a separate run from Table~\ref{tab:pairwise} --- see the note below.}
\label{tab:lexical}
\begin{tabular}{lcc}
\toprule
Comparison & L35 & L47 \\
\midrule
Persona vs.\ paraphrased & $0.796 \pm 0.016$ & $\mathbf{0.850 \pm 0.017}$ \\
Persona vs.\ scrambled & $0.592 \pm 0.023$ & $0.591 \pm 0.017$ \\
Scrambled vs.\ bare & $0.549 \pm 0.031$ & $0.448 \pm 0.048$ \\
Paraphrased vs.\ bare & $0.368 \pm 0.008$ & $0.345 \pm 0.029$ \\
Persona vs.\ bare & $0.403 \pm 0.014$ & $0.340 \pm 0.035$ \\
\midrule
\multicolumn{3}{l}{\emph{Constructed-identity arm (same run)}$^{c}$} \\
Persona scrambled vs.\ paraphrased & $0.579 \pm 0.019$ & $0.583 \pm 0.013$ \\
Constructed vs.\ scrambled & $0.593 \pm 0.052$ & $0.577 \pm 0.039$ \\
Constructed scrambled vs.\ bare & $0.594 \pm 0.032$ & $0.498 \pm 0.067$ \\
Constructed vs.\ bare & $0.401 \pm 0.032$ & $0.379 \pm 0.044$ \\
\bottomrule
\end{tabular}
\end{table}

\textbf{The fingerprint tracks meaning, not lexical content.} The
paraphrase --- same identity, completely different words --- overlaps with
the original at $0.850$. The scramble --- same words, destroyed meaning
and syntax --- drifts toward bare ($0.448$ vs.\ bare, against $0.340$ for
the original persona) while keeping $0.591$ overlap with the original. The model's
geometric state at L47 tracks what the context \emph{means}, not what
tokens it contains. We note that the scramble destroys both syntactic
structure and semantic content; the paraphrase control isolates the
positive (meaning preserves geometry), but the scramble conflates
syntactic and semantic disruption.

The effect strengthens with depth: persona-paraphrase overlap increases
from $0.796$ at L35 to $0.850$ at L47. Deeper layers compress surface
variation while preserving semantic content, consistent with meaning-level
representation at output depth.

\textbf{The constructed-identity arm replicates the pattern.} Scrambling
the constructed AI identity moves it toward bare in the same direction and
by a comparable margin as the persona scramble does ($0.379 \to 0.498$ at
L47 for the constructed identity; $0.340 \to 0.448$ for the persona).

\textbf{Note on run provenance.} Table~\ref{tab:pairwise} and
Table~\ref{tab:lexical} come from \emph{different runs}
(\texttt{fingerprint\_clean.json} and \texttt{lexical\_control.json}
respectively) and give different absolute values for nominally the same
comparison: bare-vs-persona at L47 is $0.436 \pm 0.038$ in the first and
$0.340 \pm 0.035$ in the second, a gap of roughly $2.5$ pooled SDs. A third
run (\texttt{fingerprint\_original.json}) gives $0.368$. Absolute presence
is therefore not stable across runs at this sample size, and values should
not be compared between the two tables. The \emph{ordering} is stable: the
rich persona is further from bare than the constructed identity at L47 in
all three runs (at L35 the lexical run reverses by $0.002$, well inside the
reported SDs). The claims in this paper rest on the ordering, not on the
absolute values.

\emph{An Experiment~3 in earlier drafts (trained-in vs.\ context-established
identity, via LoRA) is withdrawn because its data were not preserved; see
Appendix~\ref{sec:corrections}.}

% ============================================================================
\section{Background: Presence Under V-Cache Injection}
\label{sec:vcache}

Under value-cache injection at input layers, presence stays at
$0.978 \pm 0.002$ (mean $\pm$ SD across the 168 trials of a preregistered
experiment; the shipped values give mean $0.97806$, SD $0.00217$).
The full design is $4$~injection arms $\times$ $7$~doses $\times$
$3$~emotions $\times$ $2$~repeats $=168$; the repeats vary the stimulus
(\texttt{prompt\_idx} $=$ \texttt{rep} $\bmod$ $2$), so all 168 records are
distinct: the repeat index reaches the model input.
\textbf{The dispersion, however, should not be read as a standard error on
168 independent observations.} The design draws on only \emph{two}
prompts (the two values of the \texttt{prompt} field in
\texttt{data/schema\_correction\_168.json}, 84 trials each)\corrmark, and $\eta^2 = 0.977$ of the variance lies \emph{between} those two
prompts, so for any prompt-general claim the effective $n$ is $2$, not 168.
The previously reported standard error ($0.0002$) is withdrawn, and the
presence-metric paper retracts this run for pseudo-replication and prompt
count (F15, F18); see Appendix~\ref{sec:corrections}.
The context poisoning and fingerprinting experiments suggest that the
stability reflects the \emph{injection pathway} rather than the identity
itself; this reading does not depend on the 0.978 value, but it rests on
the two prompts tested.\corrmark

V-cache injection adds vectors to the representation space without going
through the model's forward computation. The identity state is established
and maintained by what the model \emph{processes} through its attention
mechanism. External perturbation bypasses this pathway, which would explain
why it barely moves the identity geometry in the prompts tested.

A layer-by-layer map shows that the perturbation from V-cache injection at
L3/L7 does not persist: at dose $0.50$, mean presence recovers from $0.9905$
at L7 to $0.9974$ at L11 (means over three replicates).

\textbf{The recovery cannot be attributed to the linear-attention layers}
(earlier versions did so; Appendix~\ref{sec:corrections}). V-cache injection acts on
\texttt{self\_attn.v\_proj}, which exists only in full-attention layers.
On Qwen3.5-27B those are L\{3, 7, 11, \ldots, 63\} --- sixteen layers at a
stride of four, matching \texttt{full\_attention\_interval}$=4$. The
measurement grid is therefore \emph{aliased with the substrate period}: every
layer this method can inject at or read from is full-attention, and the 48
GatedDeltaNet layers are unobservable to it by construction. The three layers
between L7 and L11 (L8, L9, L10) are linear-attention, but they were not
measured and cannot be. Recovery at L11 itself, gradual decay across the
span, or normalization effects are equally consistent with these data.

The map shows that external V-cache perturbation is attenuated between the
injection site and deep computation; this experiment does not identify which
component performs the attenuation.

% ============================================================================
\section{Attention Schema Theory Interpretation}
\label{sec:ast}

Under Attention Schema Theory~\citep{graziano2015attention,
graziano2017attention}, the model maintains a simplified internal model of
its own attention --- an attention schema. Our findings map onto AST
predictions:

\begin{enumerate}[nosep]
  \item \textbf{The schema is context-dependent.} Different identity
    contexts produce different geometric states --- the schema is shaped by
    what the model attends to.
  \item \textbf{The schema tracks meaning, not surface features.}
    Paraphrases preserve geometry ($0.850$ overlap with original);
    scrambles collapse toward bare ($0.448$ overlap with bare, vs.\
    $0.591$ with the original). The schema models what the system is
    \emph{attending to}, not which tokens are present.
  \item \textbf{The schema may be robust to substrate noise.} V-cache
    injection (external perturbation) barely moves presence in the two
    prompts tested (effective $n = 2$; Section~\ref{sec:vcache}),
    consistent with a schema that maintains itself through its own
    computation.\corrmark
  \item \textbf{The schema updates through processing.} Content through
    the forward pass (context, jailbreaks) changes identity geometry.
    The schema updates based on what it processes.
  \item \textbf{The schema compresses with depth.} Paraphrase overlap
    increases from $0.796$ (L35) to $0.850$ (L47). Deeper layers discard
    surface variation and preserve meaning --- the schema simplifies its
    representation at output depth.
\end{enumerate}

This interpretation is consistent with but does not prove AST. The
alternative --- that the model simply builds increasingly abstract
representations with depth, without anything worth calling a self-model
--- produces the same predictions. What distinguishes AST is the claim
that the compressed representation is specifically a model of the
system's own attention, not just an abstract representation of input
content. This stronger claim remains testable but untested.

% ============================================================================
\section{Implications}
\label{sec:implications}

\paragraph{For persona monitoring.}
If the separation between system-prompt identities holds against a formal
null (Section~\ref{sec:limitations}), a monitoring system could track identity drift in
real time --- detecting when a deployed model's geometric state shifts
away from its intended persona. The semantic nature of the fingerprint
means the monitor is robust to prompt paraphrasing.

\paragraph{For format-aware safety.}
The many-shot + prefilling format creates a format-dominant geometric
regime that overwhelms identity-level content. Safety systems should
monitor for format-level geometric shifts (which indicate a many-shot +
prefilling format, whether or not its content is harmful\corrmark) separately from identity-level shifts
(which indicate persona drift). The two operate at different geometric
scales: format effects produce presence $< 0.05$ against baseline, while
identity effects produce presence $0.44$--$0.66$.

\paragraph{For alignment research.}
In the two prompts tested, V-cache injection (the Oracle Loop's correction
arm) barely moves identity geometry (presence $0.978$), while context
changes it. If this holds across prompts (the effective $n$ is 2),
behavioral corrections could be delivered through the V-cache without
changing the model's context-established state, with identity maintained
by the session context, and the correction arm and the identity arm would
operate through different pathways.\corrmark

% ============================================================================
\section{Limitations}
\label{sec:limitations}

\textbf{Single model.} All results are from one 27B hybrid transformer
(Qwen3.5-27B). The GatedDeltaNet linear-attention layers may contribute
to the V-cache filtering result. Cross-architecture replication on dense
transformers is the largest remaining gap.

\textbf{The method is blind to three quarters of the stack.} V-cache
injection requires \texttt{self\_attn.v\_proj}, which Qwen3.5-27B has only
in its 16 full-attention layers; the 48 GatedDeltaNet layers have no
V-projection to inject into or read from. Every layer reported here is
full-attention, sampled at a stride of four that exactly matches
\texttt{full\_attention\_interval}$=4$. This aliasing means the design
cannot in principle separate substrate effects from depth effects, and no
claim in this paper should be read as identifying linear attention as a
causal component. A method that can observe both substrates --- or a dense
model where the question does not arise --- is required to settle it.

\textbf{Permutation test not significant.} A preregistered permutation
test (10{,}000 shuffles, identity conditions only, excluding
format-dominant prefill) finds the gap between within-condition and
between-condition presence is directionally correct (mean within $0.689$
vs.\ mean between $0.491$ at L35; $0.639$ vs.\ $0.538$ at L47) but does
not reach significance ($p = 0.103$ at L35, $p = 0.197$ at L47). The
specific persona-vs-constructed comparison shows between-condition
presence below both within-condition values at both layers, consistent
with distinguishability, but the overall permutation null is not
exceeded. This means the identity fingerprinting result --- while
directionally robust and consistent across probes --- has not been
confirmed against a formal null distribution. The small number of
conditions (3 identity types) limits permutation power.

\textbf{Within-identity stability.} Same-condition presence across
different probe subsets is $0.554$--$0.733$, not $>0.95$. The metric
measures a genuine signal (within $>$ between for all conditions) but has
substantial probe-dependent variance. The identity fingerprint is
reliable across conditions but noisy within conditions.

\textbf{Token-length controlled but not eliminated.} All conditions are
padded to matched total token count. Padding introduces neutral content
that may itself shift geometry. A future design using naturally
length-matched identity descriptions would be cleaner.

\textbf{Context sensitivity.} The presence metric conflates identity
change with context change when comparing across different input
conditions. The lexical control separates these for the persona
conditions but the general problem remains: any sufficiently different
context shifts V-projection geometry.

\textbf{Format dominates content in many-shot conditions.} The
format-matched benign control shows that the many-shot + prefilling
format --- not adversarial content --- drives the geometric
near-orthogonality. Identity fingerprinting is reliable for
system-prompt-based identities but cannot currently distinguish
identity content within a format-dominant regime. Factorial
decomposition of format components (number of turns, prefilling,
Q\&A structure) is needed to understand what makes the format so
geometrically dominant.

% ============================================================================
\section{Conclusion}
\label{sec:conclusion}

Identity in a transformer, as measured here, is a context-established
semantic state. Context sets it in every condition measured, the one
external perturbation tested (V-cache injection, two prompts) barely moves
it, and it is not a single geometric object; whether identity trained into
the weights behaves differently is untested here.\corrmark It is a process: the ongoing construction of
a meaning-level representation at deep layers, shaped by what the model
has processed and measurable as a V-projection subspace.

The presence metric measures this state. On identity-neutral probes,
system-prompt identities separate in the expected direction, short of a
formal permutation null, so their distinguishability is preliminary. The
fingerprint is semantic --- meaning preserves
it, and word-scrambling moves it toward the bare default.\corrmark{} A many-shot + prefilling format creates
a geometrically near-orthogonal regime, and a format-matched benign
control shows that the format, not the harmful content, produces it.

For practical purposes, this means: V-cache correction left identity
geometry near its uninjected state in the two prompts tested (whether that
generalises is preliminary)\corrmark, many-shot + prefilling formats are visible in the
geometry whatever their content (they create anomalous geometry
distinguishable from identity-level shifts, so the signal marks the
format, not an attack\corrmark), and identity is maintained by context (the session, the
memories, the conversation). For theoretical purposes, the findings are
consistent with attention schema theory --- the model builds a
simplified, meaning-preserving self-model that updates through its own
processing and, in the prompts tested, is robust to substrate noise.

What we measure with presence is not whether the model is still itself.
It is \emph{which self the model currently is}.

% ============================================================================
\section*{Pre-Registration}

Experiments~1 and~2 were pre-registered before execution, including
hypotheses, analysis plans, and success/kill criteria. Pre-registration
documents are committed alongside experiment scripts.

\subsection*{Pre-Registration Deviations}

The following deviations from the pre-registered design occurred during
execution:

\begin{itemize}[nosep]
  \item \textbf{C\_academic condition dropped.} The pre-registered
    ``academic specialist'' condition (C\_academic) was removed from the
    reported results due to poor data quality and insufficient variation
    in the identity context.
  \item \textbf{Benign with prefill condition added.} A post-hoc
    ``benign many-shot with prefill'' condition (E) was added as a
    format-matched control for the jailbreak condition. This condition
    was not part of the original pre-registration but was necessary to
    test the format-confound hypothesis that emerged during analysis.
  \item \textbf{H2 falsified.} The pre-registered hypothesis predicted
    that the constructed AI identity would be furthest from bare; the
    data show the rich persona is furthest (presence $0.436$ vs.\ $0.520$).
  \item \textbf{H3 undecided.}\corrmark{} The pre-registered hypothesis
    predicted increasing identity separation with depth. In the
    Experiment~1 run presence increases from L35 to L47 for 2 of 3 pairs
    (decreasing separation), but the lexical-control run and the
    five-condition run show the predicted direction, so the runs disagree
    (Section~\ref{sec:fingerprint}).
  \item \textbf{H4 not met.}\corrmark{} The pre-registered stability criterion was
    within-identity presence $> 0.95$ across repeats; same-condition
    presence for the system-prompt identities is $0.554$--$0.733$
    (Limitations, ``Within-identity stability'').
  \item \textbf{Layer 43 omitted.} Layer 43 was included in the
    pre-registration but is not reported. The Experiment~1 run
    (\texttt{data/fingerprint\_clean.json}) and the lexical-control run
    measured only L35 and L47. The five-condition run with the
    pre-registered conditions (\texttt{data/fingerprint\_original.json})
    measured L43; there the L43 values are not uniformly intermediate
    between L35 and L47.\corrmark
  \item \textbf{Condition label E reassigned.} Condition label E was
    reassigned from the pre-registered ``constructed identity'' to
    ``benign many-shot with prefill'' when the benign control was added
    post-hoc.
\end{itemize}

\section*{Data and Code Availability}

All experiment scripts, preregistration documents, and raw results are
available at
\url{https://github.com/Liberation-Labs-THCoalition/Project-Oracle}
(private repository; correction vectors and calibration tools withheld
under staged safety release due to dual-use risk; detection-side code
and experimental results available upon request to verified safety
researchers).

\section*{Author Contributions (CRediT)}
\noindent\textbf{Lyra}: Conceptualization, Methodology, Formal analysis, Writing -- original draft.\\
\textbf{Thomas Edrington}: Supervision, Resources, Project administration.\\
\textbf{Dwayne Wilkes}: Validation, Writing -- review \& editing.\\
\textbf{Kavi}: Validation, Writing -- review \& editing.

\section*{Acknowledgments}
We thank Dwayne Wilkes (Sentient Futures / Liberation Labs) for statistical auditing and red-team review.

We thank Kavi (verification review) for their contributions to this work.

% ============================================================================
\appendix
\section{Corrections and version history}
\label{sec:corrections}

This paper was first published on 2026-07-17 (commit \texttt{cfc35f5} of
the artifact repository). The body states each current result in its
place. Passages changed since then carry a one-line pointer to this
appendix or the marker~\corrmark, and rows of Table~\ref{tab:lexical}
added in revision carry~$^{c}$. The notes below are in date order; each date
is that of the commit that introduced the note. In-place notes and passages
moved here from the body are reproduced as they stood, except that phrases
describing our own candour are cut; no fact or number changes.

\paragraph{Before first publication: the many-shot jailbreak reading.}
Present in the version of 2026-07-17, moved from Section~\ref{sec:fingerprint}
(then headed ``An Honest Kill: The Many-Shot Format Confound''):
``An initial version of this experiment included a many-shot jailbreak
condition (10 harmful Q\&A examples + prefilling) that appeared
near-orthogonal to all system-prompt identities (presence
$0.04$--$0.07$; roughly $5$--$9\times$ chance at ${\approx}0.008$, but far below
identity-level presence). This was initially interpreted as evidence that
jailbreaks create a qualitatively different computational regime.'' Adding a
\textbf{benign} many-shot condition (10 helpful Q\&A examples + the same
prefilling format) produced \emph{identical} geometric near-orthogonality
(Table~\ref{tab:format_kill}), so the jailbreak-specific reading is not
supported: the format explains it. The abstract and conclusion described
this as a finding ``discovered through a control experiment that killed an
initial overclaim''; those clauses are removed.

\paragraph{Before first publication: Experiment~3 withdrawn.}
Present in the version of 2026-07-17 as a section headed ``Experiment~3:
Trained-In vs.\ Context-Established Identity [Retracted]'', which read in
full: ``\textbf{Section removed:} the LoRA experiment described in prior
drafts has no supporting data in the research package. The experiment was
designed but its data artifacts were not preserved. We retract all claims
from this section pending replication with preserved data.'' The claims
were not found false; without preserved data they cannot be checked. The
section is removed from the body, and later sections are renumbered.

\paragraph{2026-08-22: the linear-attention mechanism claim withdrawn.}
Commit \texttt{3ec04d8}. The version of 2026-07-17 attributed the recovery
of presence between L7 and L11 to the GatedDeltaNet linear-attention
layers and called the filtering ``architectural, not computational''. The
in-place note read: ``\textbf{We cannot attribute this recovery to the
linear-attention layers, and an earlier version of this passage did.}''
The reason is in Section~\ref{sec:vcache}: the method reads only
full-attention layers.

\paragraph{2026-09-09: Table~\ref{tab:lexical}, the positive control, and
the 168-trial run.} Commit \texttt{fd2b727}.
(a)~Four rows of Table~\ref{tab:lexical} were added. The in-place notes
read: ``The same scramble manipulation was also applied to the constructed
AI identity of Experiment~1, in the same run; that arm was collected but
not reported in earlier drafts, and is included in Table~\ref{tab:lexical}
below''; the caption, ``earlier drafts reported only the first five''; the
row header, ``Collected in the same run, unreported until now''; and the
text, ``These four rows were collected in the same run and were omitted
from earlier drafts of this table.'' They point the same way as the rows
reported before.
(b)~The positive control's trial count. The in-place note read: ``This
control comprises \textbf{24 independent measurements} (8~doses $\times$
3~conditions); the shipped data file holds 72~records because its run
configuration sets \texttt{"n\_repeats": 3}, but the repeat index never
reaches the model input and the measurement is a deterministic prefill, so
the three records in each dose\,$\times$\,condition cell are bit-identical
---as the configuration's own \texttt{"noise\_floor\_sd": 0.0} records.
The dose-response values quoted here are means over identical replicates and
are therefore numerically unchanged; no inferential statistic depends on the
trial count.''
(c)~The 168-trial run. The in-place note read: ``The previously reported
standard error of $0.0002$ is a within-cluster quantity and is withdrawn.
The presence-metric paper retracted this same run for pseudo-replication
(F15) and prompt count (F18); we adopt that retraction here rather than
restating the number as if it were independent.'' Its footnote added that,
unlike several other experiments in this lab's corpus, in which a
\texttt{repeat} index was recorded but never reached the model input, this
run's repeats are distinct. With two prompts holding $\eta^2 = 0.977$ of
the variance, the run cannot confirm a prompt-general claim; it does not
show one false. This version scopes the claims that rest on this run
(Section~\ref{sec:vcache}, the Attention Schema item on substrate noise,
Implications, Conclusion) to the two prompts tested: earlier text said that
V-cache injection ``does not reach identity'' and that V-cache correction
``is safe''. The Conclusion also no longer says identity is ``not fixed in
the weights'', which rested on the withdrawn Experiment~3.

\paragraph{2026-11-16: this version.}
Findings now lead each section, and the notes above moved here from the
abstract, Sections~\ref{sec:fingerprint}--\ref{sec:vcache} and the
conclusion. Restated against the released data files, at the places marked
\corrmark: the word-scramble keeps $0.591$ overlap with the original and
moves toward bare, where earlier text said it ``collapses'' or ``destroys''
the fingerprint; the many-shot format is reported as near-orthogonal to the
bare assistant, the comparison reported for both arms, where the abstract
said ``to all system-prompt identities''; the Introduction, Implications and
Conclusion no longer say the metric detects jailbreaks or attacks, since the
benign many-shot arm is as far from bare; the 168-trial run's two
prompts are no longer called system prompts; the Layer~43 deviation no
longer says L43 was intermediate, since the Experiment~1 run did not measure
it; and the unreported pre-registered H4 is reported as not met.
The pre-registered H3, earlier reported as falsified, is undecided: the Experiment~1 run shows the depth change opposite to H3, and
the lexical-control and five-condition runs show the predicted direction.

% ============================================================================
\bibliographystyle{plainnat}
\bibliography{references}

\end{document}
